\chapter{Assignment problems}
An assignment problem matches agents to tasks. Build a bipartite graph with one vertex class for agents and the other for tasks; an edge means the assignment is allowed. A matching is a set of edges with no shared endpoints. A perfect matching assigns every agent a distinct task.

For a weighted assignment, each allowed agent-task pair has a cost. The goal is to choose a perfect matching with minimum total cost. This is not the same as merely finding any perfect matching. The Hungarian algorithm solves weighted assignment problems systematically; for a small matrix, row and column reductions can help identify an optimum, but a valid solution must still assign each row and column exactly once and justify that no cheaper assignment exists. [QCE] [VCE4]

\Enrichment{} \textbf{Hall's marriage theorem.} A bipartite graph with vertex classes $U$ (agents) and $W$ (tasks) has a matching that assigns every vertex of $U$ to a different vertex of $W$ if and only if every subset $X\subseteq U$ has at least $|X|$ neighbours in $W$. Write $N(X)$ for the set of tasks adjacent to at least one agent in $X$; the condition is $|N(X)|\geq |X|$ for every $X\subseteq U$. When $|U|=|W|$, such a matching is perfect.

\begin{example}
For cost matrix $\begin{pmatrix}4&1&3\\2&0&5\\3&2&2\end{pmatrix}$, assigning row 1 to column 2, row 2 to column 1, and row 3 to column 3 has total cost $1+2+2=5$. Check the other possible permutations to determine whether this is optimal.
\end{example}

\section*{Exercises}
\begin{enumerate}
\item Use the example matrix to find the minimum assignment cost and an optimal assignment.
\item Give a bipartite graph with no perfect matching and explain the obstruction.
\item Explain why selecting the cheapest available task independently for each agent can fail.
\item \Enrichment{} Use Hall's theorem for these allowed assignments: $A$ may choose 1 or 2, $B$ may choose 1 or 2, $C$ may choose 2 or 3, and $D$ may choose 3 or 4. (a) Find a matching that assigns every agent a different task. (b) Change $C$'s choices to 1 or 2. For $X=\{A,B,C\}$, find $N(X)$ and compare $|N(X)|$ with $|X|$. (c) Explain why the modified choices cannot assign every agent a different task, even though each agent individually has options. (d) In your own words, describe what Hall's condition checks for a group of agents.
\end{enumerate}
